\input{preamble}

% OK, start here.
%
\begin{document}

\title{Catégories dérivées des champs}

\maketitle

\phantomsection
\label{section-phantom}

\tableofcontents




\section{Introduction}
\label{section-introduction}

\noindent
Dans ce chapitre, nous étudions les catégories dérivées associées aux champs
algébriques. Il s'agit notamment des catégories dérivées de faisceaux
quasi-cohérents : nous démontrons des analogues des résultats pour les schémas
(voir Catégories dérivées des schémas, section
\ref{perfect-section-introduction}) et pour les espaces algébriques (voir
Catégories dérivées des espaces, section
\ref{spaces-perfect-section-introduction}). Les résultats du présent chapitre
diffèrent de ceux de \cite{LM-B} principalement parce que nous utilisons
systématiquement les « gros sites ». Avant d'aborder ce chapitre, le lecteur
est invité à parcourir rapidement les chapitres « Faisceaux sur les champs
algébriques » et « Cohomologie des champs algébriques », où est introduite la
terminologie employée ici.



\section{Conventions, notations et abus de langage}
\label{section-conventions}

\noindent
Nous continuons d'utiliser les conventions et les abus de langage introduits
dans Propriétés des champs, section
\ref{stacks-properties-section-conventions}. Nous adoptons les notations
expliquées dans Cohomologie des champs, section
\ref{stacks-cohomology-section-notation}.












\section{Les sites lisse-\'etale et plat-fppf}
\label{section-lisse-etale}

\noindent
Cette section est l'analogue, pour les catégories dérivées, de Cohomologie des
champs, section \ref{stacks-cohomology-section-lisse-etale}.

\begin{lemma}
\label{lemma-shriek-derived}
Soit $\mathcal{X}$ un champ algébrique. Nous employons les notations de
Cohomologie des champs, lemmes
\ref{stacks-cohomology-lemma-lisse-etale} et
\ref{stacks-cohomology-lemma-lisse-etale-modules}.
\begin{enumerate}
\item Le foncteur
$g_! : \textit{Ab}(\mathcal{X}_{lisse,\etale}) \to
\textit{Ab}(\mathcal{X}_\etale)$
admet un foncteur dérivé gauche
$$
Lg_! :
D(\mathcal{X}_{lisse,\etale})
\longrightarrow
D(\mathcal{X}_\etale)
$$
qui est adjoint à gauche de $g^{-1}$ et tel que
$g^{-1}Lg_! = \text{id}$.
\item Le foncteur $g_! :
\textit{Mod}(\mathcal{X}_{lisse,\etale},
\mathcal{O}_{\mathcal{X}_{lisse,\etale}}) \to
\textit{Mod}(\mathcal{X}_\etale, \mathcal{O}_{\mathcal{X}})$
admet un foncteur dérivé gauche
$$
Lg_! :
D(\mathcal{O}_{\mathcal{X}_{lisse,\etale}})
\longrightarrow
D(\mathcal{X}_\etale, \mathcal{O}_\mathcal{X})
$$
qui est adjoint à gauche de $g^*$ et tel que $g^*Lg_! = \text{id}$.
\item Le foncteur $g_! : \textit{Ab}(\mathcal{X}_{flat,fppf}) \to
\textit{Ab}(\mathcal{X}_{fppf})$
admet un foncteur dérivé gauche
$$
Lg_! :
D(\mathcal{X}_{flat, fppf})
\longrightarrow
D(\mathcal{X}_{fppf})
$$
qui est adjoint à gauche de $g^{-1}$ et tel que
$g^{-1}Lg_! = \text{id}$.
\item Le foncteur $g_! :
\textit{Mod}(\mathcal{X}_{flat,fppf},
\mathcal{O}_{\mathcal{X}_{flat,fppf}}) \to
\textit{Mod}(\mathcal{X}_{fppf}, \mathcal{O}_{\mathcal{X}})$
admet un foncteur dérivé gauche
$$
Lg_! :
D(\mathcal{O}_{\mathcal{X}_{flat, fppf}})
\longrightarrow
D(\mathcal{O}_\mathcal{X})
$$
qui est adjoint à gauche de $g^*$ et tel que $g^*Lg_! = \text{id}$.
\end{enumerate}
Attention : il n'est pas clair a priori que $Lg_!$ sur les modules coïncide avec
$Lg_!$ sur les faisceaux abéliens ; voir Cohomologie sur les sites, remarque
\ref{sites-cohomology-remark-when-derived-shriek-equal}.
\end{lemma}

\begin{proof}
L'existence du foncteur $Lg_!$ et son adjonction à $g^*$ résultent de
Cohomologie sur les sites, lemme
\ref{sites-cohomology-lemma-existence-derived-lower-shriek}. (Dans le cas des
faisceaux abéliens, on prend le faisceau constant $\mathbf{Z}$ comme faisceau
structural.) De plus, sa valeur sur un complexe $\mathcal{H}^\bullet$ se calcule
en choisissant une résolution gauche convenable
$\mathcal{K}^\bullet \to \mathcal{H}^\bullet$, puis en appliquant le foncteur
$g_!$ à $\mathcal{K}^\bullet$. Comme
$g^{-1}g_!\mathcal{K}^\bullet = \mathcal{K}^\bullet$ d'après Cohomologie des
champs, lemmes \ref{stacks-cohomology-lemma-lisse-etale-modules} et
\ref{stacks-cohomology-lemma-lisse-etale}, la dernière assertion est vérifiée
dans chacun des cas.
\end{proof}

\begin{lemma}
\label{lemma-lisse-etale-functorial-derived}
Sous les hypothèses et avec les notations de Cohomologie des champs, lemme
\ref{stacks-cohomology-lemma-lisse-etale-functorial}, nous avons
$$
g^{-1} \circ Rf_* = Rf'_* \circ (g')^{-1}
\quad\text{et}\quad
L(g')_! \circ (f')^{-1} = f^{-1} \circ Lg_!
$$
sur les catégories dérivées non bornées, aussi bien pour les modules que pour
les faisceaux abéliens.
\end{lemma}

\begin{proof}
Soit $\tau = \etale$ (resp.\ $\tau = fppf$), et soit $\mathcal{F}$ un faisceau
abélien sur $\mathcal{X}_\tau$. D'après Cohomologie des champs, lemme
\ref{stacks-cohomology-lemma-lisse-etale-functorial-cohomology}
l'application canonique de changement de base
$$
g^{-1}Rf_*\mathcal{F} \longrightarrow Rf'_*(g')^{-1}\mathcal{F}
$$
est un isomorphisme. La suite de la démonstration est formelle. Puisque la
cohomologie des groupes abéliens coïncide avec celle des faisceaux de modules,
nous en déduisons également que
$g^{-1} Rf_*\mathcal{F} = Rf'_* (g')^{-1}\mathcal{F}$ lorsque $\mathcal{F}$ est
un faisceau de modules sur $\mathcal{X}_\tau$.

\medskip\noindent
Montrons ensuite que, pour $\mathcal{G}$, faisceau de modules ou faisceau de
groupes abéliens sur $\mathcal{Y}_{lisse,\etale}$ (resp.\
$\mathcal{Y}_{flat,fppf}$), l'application canonique
$$
L(g')_!(f')^{-1}\mathcal{G} \to f^{-1}Lg_!\mathcal{G}
$$
est un isomorphisme. Il suffit pour cela de démontrer que, pour tout faisceau
injectif $\mathcal{I}$ sur $\mathcal{X}_\tau$, l'application induite
$$
\Hom(L(g')_!(f')^{-1}\mathcal{G}, \mathcal{I}[n])
\leftarrow
\Hom(f^{-1}Lg_!\mathcal{G}, \mathcal{I}[n])
$$
est un isomorphisme pour tout $n \in \mathbf{Z}$ (les Hom étant pris dans les
catégories dérivées appropriées). Par les adjonctions entre $f^{-1}$ et $Rf_*$,
entre $Lg_!$ et $g^{-1}$, et par leurs versions « primées », cela résulte de
l'isomorphisme
$g^{-1} Rf_*\mathcal{I} \to Rf'_* (g')^{-1}\mathcal{I}$ démontré ci-dessus.

\medskip\noindent
Dans le cas d'un complexe borné $\mathcal{G}^\bullet$, de modules ou de groupes
abéliens, sur $\mathcal{Y}_{lisse,\etale}$ (resp.\
$\mathcal{Y}_{fppf}$), l'application canonique
\begin{equation}
\label{equation-to-show}
L(g')_!(f')^{-1}\mathcal{G}^\bullet \to f^{-1}Lg_!\mathcal{G}^\bullet
\end{equation}
est un isomorphisme. Cela découle du cas d'un faisceau par les arguments usuels
de troncature et du fait que les foncteurs $L(g')_!(f')^{-1}$ et
$f^{-1}Lg_!$ sont des foncteurs exacts de catégories triangulées.

\medskip\noindent
Supposons que $\mathcal{G}^\bullet$ soit un complexe borné supérieurement, de
modules ou de groupes abéliens, sur $\mathcal{Y}_{lisse,\etale}$ (resp.\
$\mathcal{Y}_{fppf}$). L'application canonique
(\ref{equation-to-show}) est un isomorphisme, car les troncatures bêtes
$\sigma_{\geq -n}$ (voir Homologie, section
\ref{homology-section-truncations}) permettent d'écrire
$\mathcal{G}^\bullet$ comme une colimite
$\mathcal{G}^\bullet = \colim \mathcal{G}_n^\bullet$
de complexes bornés. On obtient ainsi un triangle distingué
$$
\bigoplus\nolimits_{n \geq 1} \mathcal{G}_n^\bullet \to
\bigoplus\nolimits_{n \geq 1} \mathcal{G}_n^\bullet \to
\mathcal{G}^\bullet \to \ldots
$$
et chacun des foncteurs $L(g')_!$, $(f')^{-1}$, $f^{-1}$ et $Lg_!$ commute aux
sommes directes de complexes.

\medskip\noindent
Si $\mathcal{G}^\bullet$ est un complexe arbitraire, de modules ou de groupes
abéliens, sur $\mathcal{Y}_{lisse,\etale}$ (resp.\
$\mathcal{Y}_{fppf}$), nous employons les troncatures canoniques
$\tau_{\leq n}$ (voir Homologie, section
\ref{homology-section-truncations}) pour écrire $\mathcal{G}^\bullet$ comme une
colimite de complexes bornés supérieurement, puis nous reprenons l'argument du
paragraphe précédent.

\medskip\noindent
Enfin, par les adjonctions entre $f^{-1}$ et $Rf_*$, entre $Lg_!$ et $g^{-1}$,
et par leurs versions « primées », nous concluons en toute généralité que la
première identité du lemme résulte de la seconde.
\end{proof}

\begin{lemma}
\label{lemma-higher-shriek-quasi-coherent}
Soit $\mathcal{X}$ un champ algébrique. Nous employons les notations de
Cohomologie des champs, lemme
\ref{stacks-cohomology-lemma-lisse-etale}.
\begin{enumerate}
\item Soit $\mathcal{H}$ un
$\mathcal{O}_{\mathcal{X}_{lisse,\etale}}$-module quasi-cohérent sur le site
lisse-\'etale de $\mathcal{X}$. Pour tout $p \in \mathbf{Z}$, le faisceau
$H^p(Lg_!\mathcal{H})$ est un module localement quasi-cohérent sur
$\mathcal{X}$ et possède la propriété de changement de base plat.
\item Soit $\mathcal{H}$ un
$\mathcal{O}_{\mathcal{X}_{flat,fppf}}$-module quasi-cohérent sur le site
plat-fppf de $\mathcal{X}$. Pour tout $p \in \mathbf{Z}$, le faisceau
$H^p(Lg_!\mathcal{H})$ est un module localement quasi-cohérent sur
$\mathcal{X}$ et possède la propriété de changement de base plat.
\end{enumerate}
\end{lemma}

\begin{proof}
Choisissons un schéma $U$ et un morphisme lisse surjectif
$x : U \to \mathcal{X}$. D'après Modules sur les sites, définition
\ref{sites-modules-definition-site-local}, il existe un recouvrement étale
(resp.\ fppf) $\{U_i \to U\}_{i \in I}$ tel que chaque image inverse
$f_i^{-1}\mathcal{H}$ possède une présentation globale (voir Modules sur les
sites, définition \ref{sites-modules-definition-global}). Ici,
$f_i : U_i \to \mathcal{X}$ est le composé $U_i \to U \to \mathcal{X}$, qui est
un morphisme de champs algébriques. (Rappelons que l'image inverse « est » la
restriction à $\mathcal{X}/f_i$ ; voir Faisceaux sur les champs, définition
\ref{stacks-sheaves-definition-pullback} et la discussion qui suit.) Après
avoir raffiné le recouvrement, nous pouvons supposer que chaque $U_i$ est un
schéma affine. Puisque chaque $f_i$ est lisse (resp.\ plat), le lemme
\ref{lemma-lisse-etale-functorial-derived} donne
$f_i^{-1}Lg_!\mathcal{H} = Lg_{i, !}(f'_i)^{-1}\mathcal{H}$. À l'aide de
Cohomologie des champs, lemme
\ref{stacks-cohomology-lemma-check-lqc-fbc-on-covering}, nous ramenons l'énoncé
du lemme au cas où $\mathcal{H}$ possède une présentation globale et où
$\mathcal{X} = (\Sch/X)_{fppf}$ pour un schéma affine $X = \Spec(A)$.

\medskip\noindent
Écrivons cette présentation sous la forme
$$
\bigoplus\nolimits_{j \in J} \mathcal{O} \longrightarrow
\bigoplus\nolimits_{i \in I} \mathcal{O} \longrightarrow
\mathcal{H} \longrightarrow 0
$$
où $\mathcal{O} = \mathcal{O}_{\mathcal{X}_{lisse,\etale}}$ (resp.\
$\mathcal{O} = \mathcal{O}_{\mathcal{X}_{flat,fppf}}$). Remarquons que le site
$\mathcal{X}_{lisse,\etale}$ (resp.\ $\mathcal{X}_{flat,fppf}$) possède un
objet final, à savoir $X/X$, qui est quasi-compact (voir Cohomologie sur les
sites, section \ref{sites-cohomology-section-limits}). Nous avons donc
$$
\Gamma(\bigoplus\nolimits_{i \in I} \mathcal{O}) =
\bigoplus\nolimits_{i \in I} A
$$
d'après Sites, lemme \ref{sites-lemma-directed-colimits-sections}. Par
conséquent, l'application de la présentation correspond à une présentation
analogue
$$
\bigoplus\nolimits_{j \in J} A \longrightarrow
\bigoplus\nolimits_{i \in I} A \longrightarrow
M \longrightarrow 0
$$
d'un $A$-module $M$. De plus, $\mathcal{H}$ est égal à la restriction au site
lisse-\'etale (resp.\ plat-fppf) du faisceau quasi-cohérent $M^a$ associé à
$M$. Choisissons une résolution
$$
\ldots \to F_2 \to F_1 \to F_0 \to M \to 0
$$
par des $A$-modules libres. Le complexe
$$
\ldots \mathcal{O} \otimes_A F_2 \to \mathcal{O} \otimes_A F_1 \to
\mathcal{O} \otimes_A F_0 \to \mathcal{H} \to 0
$$
est une résolution de $\mathcal{H}$ par des $\mathcal{O}$-modules libres, car,
pour tout objet $U/X$ de $\mathcal{X}_{lisse,\etale}$ (resp.\
$\mathcal{X}_{flat,fppf}$), le morphisme structural $U \to X$ est plat. Par
construction, la valeur de $Lg_!\mathcal{H}$ est donc
$$
\ldots \to
\mathcal{O}_\mathcal{X} \otimes_A F_2 \to
\mathcal{O}_\mathcal{X} \otimes_A F_1 \to
\mathcal{O}_\mathcal{X} \otimes_A F_0 \to 0 \to \ldots
$$
Comme il s'agit d'un complexe de modules quasi-cohérents sur
$\mathcal{X}_\etale$ (resp.\ $\mathcal{X}_{fppf}$), il résulte de Cohomologie
des champs, proposition
\ref{stacks-cohomology-proposition-loc-qcoh-flat-base-change}, que
$H^p(Lg_!\mathcal{H})$ est quasi-cohérent.
\end{proof}




\section{Cohomologie et sites lisse-\'etale et plat-fppf}
\label{section-compare-unbounded-lisse-etale}

\noindent
Nous avons déjà vu que la cohomologie d'un faisceau sur un champ algébrique
$\mathcal{X}$ peut se calculer sur le site plat-fppf. Dans cette section, nous
démontrons qu'il en va de même pour les objets, éventuellement non bornés, de
la catégorie dérivée de $\mathcal{X}$.

\begin{lemma}
\label{lemma-higher-shriek-Z}
Soit $\mathcal{X}$ un champ algébrique. Nous avons
$Lg_!\mathbf{Z} = \mathbf{Z}$, qu'il s'agisse du $Lg_!$ de la partie (1) du
lemme \ref{lemma-shriek-derived} ou du $Lg_!$ de la partie (3) du lemme
\ref{lemma-shriek-derived}.
\end{lemma}

\begin{proof}
Nous le démontrons pour la comparaison du site plat-fppf au site fppf ; le cas
du site lisse-\'etale est exactement analogue. Il faut montrer que
$H^i(Lg_!\mathbf{Z})$ est $0$ pour $i \not = 0$ et que l'application canonique
$H^0(Lg_!\mathbf{Z}) \to \mathbf{Z}$ est un isomorphisme. Soit
$f : \mathcal{U} \to \mathcal{X}$ un morphisme plat et surjectif, où
$\mathcal{U}$ est un schéma, et supposons en outre $f$ localement de
présentation finie. (On peut, par exemple, choisir une présentation
$U \to \mathcal{X}$ et prendre pour $\mathcal{U}$ le champ algébrique
correspondant à $U$.) D'après Faisceaux sur les champs, lemmes
\ref{stacks-sheaves-lemma-check-exactness-covering} et
\ref{stacks-sheaves-lemma-surjective-flat-locally-finite-presentation}
il suffit de montrer que l'image inverse $f^{-1}H^i(Lg_!\mathbf{Z})$ est $0$
pour $i \not = 0$ et que l'image inverse
$H^0(Lg_!\mathbf{Z}) \to f^{-1}\mathbf{Z}$ est un isomorphisme. Le lemme
\ref{lemma-lisse-etale-functorial-derived} donne
$f^{-1}Lg_!\mathbf{Z} = L(g')_!\mathbf{Z}$, où
$g' : \Sh(\mathcal{U}_{flat, fppf}) \to \Sh(\mathcal{U}_{fppf})$
est le morphisme de comparaison correspondant pour $\mathcal{U}$. Nous sommes
ainsi ramenés au cas étudié au paragraphe suivant.

\medskip\noindent
Supposons $\mathcal{X} = (\Sch/X)_{fppf}$ pour un schéma $X$. La catégorie
$\mathcal{X}_{flat, fppf}$ possède alors un objet final $e$, à savoir $X/X$ ;
de plus, le foncteur
$u : \mathcal{X}_{flat, fppf} \to \mathcal{X}_{fppf}$ envoie $e$ sur l'objet
final. Puisque $\mathbf{Z}$ est le faisceau abélien libre sur l'objet final,
lorsque celui-ci existe, nous obtenons $Lg_!\mathbf{Z} = \mathbf{Z}$ par la
construction même de $Lg_!$ dans Cohomologie sur les sites, lemme
\ref{sites-cohomology-lemma-existence-derived-lower-shriek}.
\end{proof}

\begin{lemma}
\label{lemma-lisse-etale-cohomology}
Soit $\mathcal{X}$ un champ algébrique. Nous employons les notations du lemme
\ref{lemma-shriek-derived}.
\begin{enumerate}
\item Pour $K$ dans $D(\mathcal{X}_\etale)$, nous avons
\begin{enumerate}
\item $R\Gamma(\mathcal{X}_\etale, K) =
R\Gamma(\mathcal{X}_{lisse,\etale}, g^{-1}K)$, et
\item $R\Gamma(x, K) =
R\Gamma(\mathcal{X}_{lisse,\etale}/x, g^{-1}K)$
pour tout objet $x$ de $\mathcal{X}_{lisse,\etale}$.
\end{enumerate}
\item Pour $K$ dans $D(\mathcal{X}_{fppf})$, nous avons
\begin{enumerate}
\item $R\Gamma(\mathcal{X}_{fppf}, K) =
R\Gamma(\mathcal{X}_{flat,fppf}, g^{-1}K)$, et
\item $H^p(x, K) =
R\Gamma(\mathcal{X}_{flat,fppf}/x, g^{-1}K)$
pour tout objet $x$ de $\mathcal{X}_{flat,fppf}$.
\end{enumerate}
\end{enumerate}
Dans les deux cas, les mêmes assertions valent pour les modules, puisque
$g^{-1} = g^*$ et que le calcul de la cohomologie ne présente aucune différence
d'après Cohomologie sur les sites, lemme
\ref{sites-cohomology-lemma-modules-abelian-unbounded}.
\end{lemma}

\begin{proof}
Nous le démontrons pour la comparaison du site plat-fppf au site fppf ; le cas
du site lisse-\'etale est exactement analogue. D'après le lemme
\ref{lemma-higher-shriek-Z}, nous avons $Lg_!\mathbf{Z} = \mathbf{Z}$. Il vient
\begin{align*}
R\Gamma(\mathcal{X}_{fppf}, K)
& =
R\Hom(\mathbf{Z}, K) \\
& =
R\Hom(Lg_!\mathbf{Z}, K) \\
& =
R\Hom(\mathbf{Z}, g^{-1}K) \\
& =
R\Gamma(\mathcal{X}_{lisse,\etale}, g^{-1}K)
\end{align*}
Ceci démontre (1)(a). La partie (1)(b) résulte de la partie (1)(a). En effet,
si $x$ est au-dessus du schéma $U$, alors le site $\mathcal{X}_\etale/x$ est
équivalent à $(\Sch/U)_\etale$, et $\mathcal{X}_{lisse,\etale}$ est équivalent
à $U_{lisse, \etale}$.
\end{proof}










\section{Catégories dérivées de modules quasi-cohérents}
\label{section-derived}

\noindent
Soit $\mathcal{X}$ un champ algébrique. Comme le foncteur d'inclusion
$\QCoh(\mathcal{O}_\mathcal{X}) \to
\textit{Mod}(\mathcal{O}_\mathcal{X})$ n'est pas exact, nous ne pouvons pas
définir $D_\QCoh(\mathcal{O}_\mathcal{X})$ comme la sous-catégorie pleine de
$D(\mathcal{O}_\mathcal{X})$ formée des complexes dont les faisceaux de
cohomologie sont quasi-cohérents. Nous définissons plutôt la catégorie dérivée
des modules quasi-cohérents comme un quotient, par analogie avec Cohomologie
des champs, remarque
\ref{stacks-cohomology-remark-bousfield-colocalization}.

\medskip\noindent
Rappelons que $\textit{LQCoh}^{fbc}(\mathcal{O}_\mathcal{X}) \subset
\textit{Mod}(\mathcal{O}_\mathcal{X})$ désigne la sous-catégorie pleine des
$\mathcal{O}_\mathcal{X}$-modules localement quasi-cohérents possédant la
propriété de changement de base plat ; voir Cohomologie des champs, section
\ref{stacks-cohomology-section-loc-qcoh-flat-base-change}.
Nous abrégerons
$$
D_{\textit{LQCoh}^{fbc}}(\mathcal{O}_\mathcal{X}) =
D_{\textit{LQCoh}^{fbc}(\mathcal{O}_\mathcal{X})}(\mathcal{O}_\mathcal{X})
$$
Du lemme \ref{derived-lemma-cohomology-in-serre-subcategory} de Catégories
dérivées et de la partie (2) de Cohomologie des champs, proposition
\ref{stacks-cohomology-proposition-loc-qcoh-flat-base-change}, nous déduisons
que $D_{\textit{LQCoh}^{fbc}}(\mathcal{O}_\mathcal{X})$ est une sous-catégorie
triangulée strictement pleine et saturée de
$D(\mathcal{O}_\mathcal{X})$.

\medskip\noindent
Notons $\textit{Parasitic}(\mathcal{O}_\mathcal{X}) \subset
\textit{Mod}(\mathcal{O}_\mathcal{X})$ la sous-catégorie pleine des
$\mathcal{O}_\mathcal{X}$-modules parasites ; voir Cohomologie des champs,
section
\ref{stacks-cohomology-section-parasitic}.
Abrégeons de même
$$
D_{\textit{Parasitic}}(\mathcal{O}_\mathcal{X}) =
D_{\textit{Parasitic}(\mathcal{O}_\mathcal{X})}(\mathcal{O}_\mathcal{X})
$$
Comme précédemment, c'est une sous-catégorie triangulée strictement pleine et
saturée de $D(\mathcal{O}_\mathcal{X})$, puisque
$\textit{Parasitic}(\mathcal{O}_\mathcal{X})$ est une sous-catégorie de Serre
de $\textit{Mod}(\mathcal{O}_\mathcal{X})$ ; voir Cohomologie des champs, lemme
\ref{stacks-cohomology-lemma-parasitic}.

\medskip\noindent
L'intersection des sous-catégories faiblement de Serre
$\textit{Parasitic}(\mathcal{O}_\mathcal{X}) \cap
\textit{LQCoh}^{fbc}(\mathcal{O}_\mathcal{X})$
de $\textit{Mod}(\mathcal{O}_\mathcal{X})$ est encore une sous-catégorie
faiblement de Serre. Nous adoptons de même l'abréviation
\begin{align*}
D_{\textit{Parasitic} \cap \textit{LQCoh}^{fbc}}(\mathcal{O}_\mathcal{X})
& =
D_{\textit{Parasitic}(\mathcal{O}_\mathcal{X}) \cap
\textit{LQCoh}^{fbc}(\mathcal{O}_\mathcal{X})}(\mathcal{O}_\mathcal{X}) \\
& =
D_{\textit{Parasitic}}(\mathcal{O}_\mathcal{X})
\cap
D_{\textit{LQCoh}^{fbc}}(\mathcal{O}_\mathcal{X})
\end{align*}
Comme précédemment, c'est une sous-catégorie triangulée strictement pleine et
saturée de $D(\mathcal{O}_\mathcal{X})$. A fortiori, c'est donc une
sous-catégorie triangulée strictement pleine et saturée de
$D_{\textit{LQCoh}^{fbc}}(\mathcal{O}_\mathcal{X})$.

\begin{definition}
\label{definition-derived}
Soit $\mathcal{X}$ un champ algébrique. Avec les notations ci-dessus, nous
définissons la {\it catégorie dérivée des $\mathcal{O}_\mathcal{X}$-modules à
faisceaux de cohomologie quasi-cohérents} comme le quotient de
Verdier\footnote{Cette définition diffère de celle de la littérature ; voir
\cite[6.3]{olsson_sheaves}. Elle coïncide toutefois avec cette dernière d'après
le lemme \ref{lemma-derived-quasi-coherent}.}
$$
D_\QCoh(\mathcal{O}_\mathcal{X}) =
D_{\textit{LQCoh}^{fbc}}(\mathcal{O}_\mathcal{X})/
D_{\textit{Parasitic} \cap \textit{LQCoh}^{fbc}}(\mathcal{O}_\mathcal{X})
$$
\end{definition}

\noindent
Le quotient de Verdier est défini dans Catégories dérivées, section
\ref{derived-section-quotients}. Un morphisme $a : E \to E'$ de
$D_{\textit{LQCoh}^{fbc}}(\mathcal{O}_\mathcal{X})$ devient un isomorphisme
dans $D_\QCoh(\mathcal{O}_\mathcal{X})$ si et seulement si le cône $C(a)$ a des
faisceaux de cohomologie parasites ; voir Catégories dérivées, lemme
\ref{derived-lemma-operations}.

\medskip\noindent
Considérons les foncteurs
$$
D_{\textit{LQCoh}^{fbc}}(\mathcal{O}_\mathcal{X})
\xrightarrow{H^i}
\textit{LQCoh}^{fbc}(\mathcal{O}_\mathcal{X})
\xrightarrow{Q}
\QCoh(\mathcal{O}_\mathcal{X})
$$
Remarquons que $Q$ annule la sous-catégorie
$\textit{Parasitic}(\mathcal{O}_\mathcal{X}) \cap
\textit{LQCoh}^{fbc}(\mathcal{O}_\mathcal{X})$ ; voir Cohomologie des champs,
lemme \ref{stacks-cohomology-lemma-adjoint-kernel-parasitic}. D'après
Catégories dérivées, lemme
\ref{derived-lemma-universal-property-quotient}, nous obtenons un foncteur
cohomologique
\begin{equation}
\label{equation-Hi-quasi-coherent}
H^i :
D_\QCoh(\mathcal{O}_\mathcal{X})
\longrightarrow
\QCoh(\mathcal{O}_\mathcal{X})
\end{equation}
De plus, remarquons que $E \in D_\QCoh(\mathcal{O}_\mathcal{X})$ est nul si et
seulement si $H^i(E) = 0$ pour tout $i \in \mathbf{Z}$, puisque le noyau de
$Q$ est exactement
$\textit{Parasitic}(\mathcal{O}_\mathcal{X}) \cap
\textit{LQCoh}^{fbc}(\mathcal{O}_\mathcal{X})$ d'après Cohomologie des champs,
lemme \ref{stacks-cohomology-lemma-adjoint-kernel-parasitic}.

\medskip\noindent
Remarquons que les catégories
$\textit{Parasitic}(\mathcal{O}_\mathcal{X}) \cap
\textit{LQCoh}^{fbc}(\mathcal{O}_\mathcal{X})$ et
$\textit{LQCoh}^{fbc}(\mathcal{O}_\mathcal{X})$
constituent aussi des sous-catégories faiblement de Serre de la catégorie
abélienne
$\textit{Mod}(\mathcal{X}_\etale, \mathcal{O}_\mathcal{X})$
des modules pour la topologie étale ; voir Cohomologie des champs, proposition
\ref{stacks-cohomology-proposition-loc-qcoh-flat-base-change} et lemme
\ref{stacks-cohomology-lemma-parasitic}. L'énoncé du lemme suivant a donc un
sens.

\begin{lemma}
\label{lemma-compare-etale-fppf-QCoh}
Soit $\mathcal{X}$ un champ algébrique. Posons, par abréviation,
$\mathcal{P}_\mathcal{X} = \textit{Parasitic}(\mathcal{O}_\mathcal{X}) \cap
\textit{LQCoh}^{fbc}(\mathcal{O}_\mathcal{X})$.
Le morphisme de comparaison
$\epsilon : \mathcal{X}_{fppf} \to \mathcal{X}_\etale$
induit un diagramme commutatif
$$
\xymatrix{
D_{\textit{Parasitic} \cap \textit{LQCoh}^{fbc}}(\mathcal{O}_\mathcal{X})
\ar[r] &
D_{\textit{LQCoh}^{fbc}}(\mathcal{O}_\mathcal{X}) \ar[r] &
D(\mathcal{O}_\mathcal{X}) \\
D_{\mathcal{P}_\mathcal{X}}(\mathcal{X}_\etale, \mathcal{O}_\mathcal{X})
\ar[r] \ar[u]^{\epsilon^*} &
D_{\textit{LQCoh}^{fbc}(\mathcal{O}_\mathcal{X})}(
\mathcal{X}_\etale, \mathcal{O}_\mathcal{X})
\ar[r] \ar[u]^{\epsilon^*} &
D(\mathcal{X}_\etale, \mathcal{O}_\mathcal{X})
\ar[u]^{\epsilon^*}
}
$$
De plus, les deux flèches verticales de gauche sont des équivalences de
catégories triangulées ; nous obtenons donc aussi une équivalence
$$
D_{\textit{LQCoh}^{fbc}(\mathcal{O}_\mathcal{X})}
(\mathcal{X}_\etale, \mathcal{O}_\mathcal{X})
/
D_{\mathcal{P}_\mathcal{X}}(\mathcal{X}_\etale, \mathcal{O}_\mathcal{X})
\longrightarrow
D_\QCoh(\mathcal{O}_\mathcal{X})
$$
\end{lemma}

\begin{proof}
Puisque $\epsilon^*$ est exact, il est clair que nous obtenons le diagramme de
l'énoncé. Montrons que la flèche verticale médiane est une équivalence en
appliquant Cohomologie sur les sites, lemme
\ref{sites-cohomology-lemma-compare-topologies-derived-adequate-modules}
à la situation suivante :
$\mathcal{C} = \mathcal{X}$,
$\tau = fppf$,
$\tau' = \etale$,
$\mathcal{O} = \mathcal{O}_\mathcal{X}$,
$\mathcal{A} = \textit{LQCoh}^{fbc}(\mathcal{O}_\mathcal{X})$, et
$\mathcal{B}$ est l'ensemble des objets de $\mathcal{X}$ situés au-dessus de
schémas affines. Pour appliquer le lemme, il faut vérifier les conditions (1),
(2), (3) et (4). Les conditions (1) et (2) résultent clairement de la
discussion précédente ; plus précisément, elles découlent de Cohomologie des
champs, proposition
\ref{stacks-cohomology-proposition-loc-qcoh-flat-base-change}. La condition
(3) est satisfaite parce que tout schéma possède un recouvrement ouvert de
Zariski par des schémas affines. La condition (4) résulte de Descente, lemme
\ref{descent-lemma-quasi-coherent-and-flat-base-change}.

\medskip\noindent
Nous omettons de vérifier que l'équivalence de catégories $\epsilon^* :
D_{\textit{LQCoh}^{fbc}(\mathcal{O}_\mathcal{X})}(
\mathcal{X}_\etale, \mathcal{O}_\mathcal{X})
\to
D_{\textit{LQCoh}^{fbc}}(\mathcal{O}_\mathcal{X})$
induit une équivalence entre les sous-catégories de complexes à faisceaux de
cohomologie parasites.
\end{proof}

\noindent
Soit $\mathcal{X}$ un champ algébrique. D'après Cohomologie des champs, lemme
\ref{stacks-cohomology-lemma-quasi-coherent-weak-serre}
la catégorie des modules quasi-cohérents
$\QCoh(\mathcal{O}_{\mathcal{X}_{lisse,\etale}})$
forme une sous-catégorie faiblement de Serre de
$\textit{Mod}(\mathcal{O}_{\mathcal{X}_{lisse,\etale}})$
et la catégorie des modules quasi-cohérents
$\QCoh(\mathcal{O}_{\mathcal{X}_{flat,fppf}})$
forme une sous-catégorie faiblement de Serre de
$\textit{Mod}(\mathcal{O}_{\mathcal{X}_{flat,fppf}})$.
Nous pouvons donc considérer
$$
D_\QCoh(\mathcal{O}_{\mathcal{X}_{lisse,\etale}}) =
D_{\QCoh(\mathcal{O}_{\mathcal{X}_{lisse,\etale}})}(
\mathcal{O}_{\mathcal{X}_{lisse,\etale}})
\subset
D(\mathcal{O}_{\mathcal{X}_{lisse,\etale}})
$$
et, de même,
$$
D_\QCoh(\mathcal{O}_{\mathcal{X}_{flat,fppf}}) =
D_{\QCoh(\mathcal{O}_{\mathcal{X}_{flat,fppf}})}(
\mathcal{O}_{\mathcal{X}_{flat,fppf}})
\subset
D(\mathcal{O}_{\mathcal{X}_{flat,fppf}})
$$
Comme ci-dessus, ce sont des sous-catégories triangulées strictement pleines et
saturées. Il se trouve que $D_\QCoh(\mathcal{O}_\mathcal{X})$ est équivalente à
chacune d'elles.

\begin{lemma}
\label{lemma-derived-quasi-coherent}
Soit $\mathcal{X}$ un champ algébrique. Posons
$\mathcal{P}_\mathcal{X} = \textit{Parasitic}(\mathcal{O}_\mathcal{X}) \cap
\textit{LQCoh}^{fbc}(\mathcal{O}_\mathcal{X})$.
\begin{enumerate}
\item
Soit $\mathcal{F}^\bullet$ un objet de
$D_{\textit{LQCoh}^{fbc}(\mathcal{O}_\mathcal{X})}
(\mathcal{X}_\etale, \mathcal{O}_\mathcal{X})$.
Pour le site lisse-\'etale, prenons $g$ comme dans Cohomologie des champs,
lemme \ref{stacks-cohomology-lemma-lisse-etale}. Nous avons
\begin{enumerate}
\item $g^*\mathcal{F}^\bullet$ appartient à
$D_\QCoh(\mathcal{O}_{\mathcal{X}_{lisse,\etale}})$,
\item $g^*\mathcal{F}^\bullet = 0$ si et seulement si
$\mathcal{F}^\bullet$ appartient à
$D_{\mathcal{P}_\mathcal{X}}(\mathcal{X}_\etale, \mathcal{O}_\mathcal{X})$,
\item $Lg_!\mathcal{H}^\bullet$ appartient à
$D_{\textit{LQCoh}^{fbc}(\mathcal{O}_\mathcal{X})}(
\mathcal{X}_\etale, \mathcal{O}_\mathcal{X})$
pour $\mathcal{H}^\bullet$ appartenant à
$D_\QCoh(\mathcal{O}_{\mathcal{X}_{lisse,\etale}})$, et
\item les foncteurs $g^*$ et $Lg_!$ définissent des foncteurs mutuellement
inverses
$$
\xymatrix{
D_\QCoh(\mathcal{O}_\mathcal{X}) \ar@<1ex>[r]^-{g^*} &
D_\QCoh(\mathcal{O}_{\mathcal{X}_{lisse,\etale}})
\ar@<1ex>[l]^-{Lg_!}
}
$$
\end{enumerate}
\item
Soit $\mathcal{F}^\bullet$ un objet de
$D_{\textit{LQCoh}^{fbc}}(\mathcal{O}_\mathcal{X})$. Pour le site plat-fppf,
prenons $g$ comme dans Cohomologie des champs, lemme
\ref{stacks-cohomology-lemma-lisse-etale}. Nous avons
\begin{enumerate}
\item $g^*\mathcal{F}^\bullet$ appartient à
$D_\QCoh(\mathcal{O}_{\mathcal{X}_{flat, fppf}})$,
\item $g^*\mathcal{F}^\bullet = 0$ si et seulement si
$\mathcal{F}^\bullet$ appartient à
$D_{\mathcal{P}_\mathcal{X}}(\mathcal{O}_\mathcal{X})$,
\item $Lg_!\mathcal{H}^\bullet$ appartient à
$D_{\textit{LQCoh}^{fbc}}(\mathcal{O}_\mathcal{X})$
pour $\mathcal{H}^\bullet$ appartenant à
$D_\QCoh(\mathcal{O}_{\mathcal{X}_{flat,fppf}})$, et
\item les foncteurs $g^*$ et $Lg_!$ définissent des foncteurs mutuellement
inverses
$$
\xymatrix{
D_\QCoh(\mathcal{O}_\mathcal{X}) \ar@<1ex>[r]^-{g^*} &
D_\QCoh(\mathcal{O}_{\mathcal{X}_{flat,fppf}}) \ar@<1ex>[l]^-{Lg_!}
}
$$
\end{enumerate}
\end{enumerate}
\end{lemma}

\begin{proof}
Le foncteur $g^* = g^{-1}$ est exact ; (1)(a), (2)(a), (1)(b) et (2)(b)
résultent donc de Cohomologie des champs, lemmes
\ref{stacks-cohomology-lemma-quasi-coherent} et
\ref{stacks-cohomology-lemma-parasitic-in-terms-flat-fppf}.

\medskip\noindent
Démonstration de (1)(c) et (2)(c). La construction de $Lg_!$ au lemme
\ref{lemma-shriek-derived} (par l'intermédiaire de Cohomologie sur les sites,
lemme \ref{sites-cohomology-lemma-existence-derived-lower-shriek}, qui utilise
à son tour Catégories dérivées, proposition
\ref{derived-proposition-left-derived-exists}) montre que $Lg_!$, sur tout objet
$\mathcal{H}^\bullet$ de
$D(\mathcal{O}_{\mathcal{X}_{lisse,\etale}})$, se calcule par
$$
Lg_!\mathcal{H}^\bullet = \colim g_!\mathcal{K}_n^\bullet =
g_! \colim \mathcal{K}_n^\bullet
$$
(colimites terme à terme), où le quasi-isomorphisme
$\colim \mathcal{K}_n^\bullet \to \mathcal{H}^\bullet$ induit des
quasi-isomorphismes
$\mathcal{K}_n^\bullet \to \tau_{\leq n} \mathcal{H}^\bullet$. Comme les
foncteurs d'inclusion
$$
\textit{LQCoh}^{fbc}(\mathcal{O}_\mathcal{X}) \subset
\textit{Mod}(\mathcal{X}_\etale, \mathcal{O}_\mathcal{X})
\quad\text{et}\quad
\textit{LQCoh}^{fbc}(\mathcal{O}_\mathcal{X}) \subset
\textit{Mod}(\mathcal{O}_\mathcal{X})
$$
sont compatibles aux colimites filtrantes, il suffit de démontrer (c) pour les
complexes bornés supérieurement $\mathcal{H}^\bullet$ appartenant à
$D_\QCoh(\mathcal{O}_{\mathcal{X}_{lisse,\etale}})$ et à
$D_\QCoh(\mathcal{O}_{\mathcal{X}_{flat,fppf}})$. Dans ce cas, pour montrer que
$H^n(Lg_!\mathcal{H}^\bullet)$ appartient à
$\textit{LQCoh}^{fbc}(\mathcal{O}_\mathcal{X})$, nous pouvons raisonner par
récurrence sur l'entier $m$ tel que $\mathcal{H}^i = 0$ pour $i > m$. Si
$m < n$, alors $H^n(Lg_!\mathcal{H}^\bullet) = 0$ et le résultat est établi.
Dans le cas général, considérons le triangle distingué
$$
\tau_{\leq m - 1}\mathcal{H}^\bullet \to \mathcal{H}^\bullet \to
H^m(\mathcal{H}^\bullet)[-m] \to \ldots
$$
(Catégories dérivées, remarque
\ref{derived-remark-truncation-distinguished-triangle})
et appliquons le foncteur $Lg_!$. Comme
$\textit{LQCoh}^{fbc}(\mathcal{O}_\mathcal{X})$
est une sous-catégorie faiblement de Serre de la catégorie des modules, il
suffit de démontrer (c) pour deux termes sur trois. Le résultat vaut pour
$Lg_!\tau_{\leq m - 1}\mathcal{H}^\bullet$ par récurrence et pour
$Lg_!H^m(\mathcal{H}^\bullet)[-m]$ d'après le lemme
\ref{lemma-higher-shriek-quasi-coherent}. Par conséquent, (c) est vérifié.

\medskip\noindent
Démontrons (2)(d). D'après (2)(a) et (2)(b), le foncteur $g^{-1} = g^*$ induit
un foncteur
$$
c :
D_\QCoh(\mathcal{O}_\mathcal{X})
\longrightarrow
D_\QCoh(\mathcal{O}_{\mathcal{X}_{flat, fppf}})
$$
voir Catégories dérivées, lemme
\ref{derived-lemma-universal-property-quotient}. Nous avons donc le diagramme
suivant de catégories triangulées
$$
\xymatrix{
D_{\textit{LQCoh}^{fbc}}(\mathcal{O}_\mathcal{X})
\ar[rd]^{g^{-1}} \ar[rr]_q & &
D_\QCoh(\mathcal{O}_\mathcal{X}) \ar[ld]^c \\
& D_\QCoh(\mathcal{O}_{\mathcal{X}_{flat, fppf}})
\ar@<1ex>[lu]^{Lg_!}
}
$$
où $q$ est le foncteur quotient, le triangle intérieur est commutatif et
$g^{-1}Lg_! = \text{id}$. Pour tout objet $E$ de
$D_{\textit{LQCoh}^{fbc}}(\mathcal{O}_\mathcal{X})$, l'image de l'application
$a : Lg_!g^{-1}E \to E$ est un quasi-isomorphisme dans
$D(\mathcal{O}_{\mathcal{X}_{flat, fppf}})$. Le cône de $a$ est donc envoyé
sur zéro par $g^{-1}$ et, d'après (2)(b), $q(a)$ est un isomorphisme. Ainsi,
$q \circ Lg_!$ est un quasi-inverse de $c$.

\medskip\noindent
Dans le cas du site lisse-\'etale, le même argument démontre que
$$
D_{\textit{LQCoh}^{fbc}(\mathcal{O}_\mathcal{X})}(
\mathcal{X}_\etale, \mathcal{O}_\mathcal{X})
/
D_{\mathcal{P}_\mathcal{X}}(
\mathcal{X}_\etale, \mathcal{O}_\mathcal{X})
$$
est équivalent à
$D_\QCoh(\mathcal{O}_{\mathcal{X}_{lisse,\etale}})$. L'application de la
dernière équivalence du lemme \ref{lemma-compare-etale-fppf-QCoh} achève la
démonstration.
\end{proof}

\noindent
Le lemme suivant montre que le foncteur quotient
$D_{\textit{LQCoh}^{fbc}}(\mathcal{O}_\mathcal{X}) \to
D_\QCoh(\mathcal{O}_\mathcal{X})$ admet un adjoint à gauche. Voir la remarque
\ref{remark-QCoh-admissible}.

\begin{lemma}
\label{lemma-bousfield-colocalization}
Soit $\mathcal{X}$ un champ algébrique. Soit $E$ un objet de
$D_{\textit{LQCoh}^{fbc}}(\mathcal{O}_\mathcal{X})$. Il existe un triangle
distingué canonique
$$
E' \to E \to P \to E'[1]
$$
dans $D_{\textit{LQCoh}^{fbc}}(\mathcal{O}_\mathcal{X})$, tel que $P$
appartienne à $D_{\textit{Parasitic} \cap \textit{LQCoh}^{fbc}}
(\mathcal{O}_\mathcal{X})$
et que
$$
\Hom_{D(\mathcal{O}_\mathcal{X})}(E', P') = 0
$$
pour tout $P'$ appartenant à
$D_{\textit{Parasitic} \cap \textit{LQCoh}^{fbc}}(\mathcal{O}_\mathcal{X})$.
\end{lemma}

\begin{proof}
Considérons le morphisme de topos annelés
$g : \Sh(\mathcal{X}_{flat, fppf}) \longrightarrow \Sh(\mathcal{X}_{fppf})$
étudié dans Cohomologie des champs, section
\ref{stacks-cohomology-section-lisse-etale}.
Posons $E' = Lg_!g^*E$ et soit $P$ le cône de l'application d'adjonction
$E' \to E$ ; voir la partie (4) du lemme \ref{lemma-shriek-derived}. D'après
les parties (2)(a) et (2)(c) du lemme
\ref{lemma-derived-quasi-coherent}, $E'$ appartient à
$D_{\textit{LQCoh}^{fbc}}(\mathcal{O}_\mathcal{X})$. Par conséquent, $P$
appartient aussi à $D_{\textit{LQCoh}^{fbc}}(\mathcal{O}_\mathcal{X})$.
L'application $g^*E' \to g^*E$ est un isomorphisme, car
$g^*Lg_! = \text{id}$ d'après la partie (4) du lemme
\ref{lemma-shriek-derived}. Ainsi, $g^*P = 0$, de sorte que $P$ est un objet de
$D_{\textit{Parasitic} \cap \textit{LQCoh}^{fbc}}(\mathcal{O}_\mathcal{X})$
d'après la partie (2)(b) du lemme \ref{lemma-derived-quasi-coherent}. Enfin,
pour $P'$ appartenant à
$D_{\textit{Parasitic} \cap \textit{LQCoh}^{fbc}}(\mathcal{O}_\mathcal{X})$
nous avons
$$
\Hom(E', P') = \Hom(Lg_!g^*E, P') = \Hom(g^*E, g^*P') = 0
$$
car $g^*P' = 0$ d'après la partie (2)(b) du lemme
\ref{lemma-derived-quasi-coherent}. Le triangle distingué
$E' \to E \to P \to E'[1]$ est canonique — plus précisément, il est unique à
un isomorphisme de triangles induisant l'identité sur $E$ près — d'après la
discussion de Catégories dérivées, section \ref{derived-section-admissible}.
\end{proof}

\begin{remark}
\label{remark-QCoh-admissible}
Le lemme \ref{lemma-bousfield-colocalization} montre que
$$
D_{\textit{Parasitic} \cap \textit{LQCoh}^{fbc}}(\mathcal{O}_\mathcal{X})
\subset
D_{\textit{LQCoh}^{fbc}}(\mathcal{O}_\mathcal{X})
$$
est une sous-catégorie admissible à gauche ; voir Catégories dérivées, section
\ref{derived-section-admissible}. En particulier, si
$\mathcal{A} \subset D_{\textit{LQCoh}^{fbc}}(\mathcal{O}_\mathcal{X})$
désigne son orthogonal gauche, alors Catégories dérivées, proposition
\ref{derived-proposition-summarize-admissible}, implique que $\mathcal{A}$ est
admissible à droite dans
$D_{\textit{LQCoh}^{fbc}}(\mathcal{O}_\mathcal{X})$ et que le composé
$$
\mathcal{A} \longrightarrow
D_{\textit{LQCoh}^{fbc}}(\mathcal{O}_\mathcal{X}) \longrightarrow
D_\QCoh(\mathcal{O}_\mathcal{X})
$$
est une équivalence. Cela signifie que nous pouvons considérer
$D_\QCoh(\mathcal{O}_\mathcal{X})$ comme une sous-catégorie triangulée
strictement pleine et saturée de
$D_{\textit{LQCoh}^{fbc}}(\mathcal{O}_\mathcal{X})$, ainsi que de
$D(\mathcal{X}_{fppf}, \mathcal{O}_\mathcal{X})$.
\end{remark}





\section{Image directe dérivée des modules quasi-cohérents}
\label{section-derived-pushforward}

\noindent
Comme première application de ce qui précède, nous construisons l'image directe
dérivée. Dans Exemples, section
\ref{examples-section-derived-push-quasi-coherent}, le lecteur trouvera un
exemple de morphisme quasi-compact et quasi-séparé
$f : \mathcal{X} \to \mathcal{Y}$ de champs algébriques tel que le foncteur
image directe $Rf_*$ n'induise pas de foncteur
$D_\QCoh(\mathcal{O}_\mathcal{X}) \to
D_\QCoh(\mathcal{O}_\mathcal{Y})$. Il est donc nécessaire de se restreindre aux
complexes bornés inférieurement.

\begin{proposition}
\label{proposition-derived-direct-image-quasi-coherent}
Soit $f : \mathcal{X} \to \mathcal{Y}$ un morphisme quasi-compact et
quasi-séparé de champs algébriques. Le foncteur $Rf_*$ induit un diagramme
commutatif
$$
\xymatrix{
D^{+}_{\textit{Parasitic} \cap \textit{LQCoh}^{fbc}}(\mathcal{O}_\mathcal{X})
\ar[r] \ar[d]^{Rf_*} &
D^{+}_{\textit{LQCoh}^{fbc}}(\mathcal{O}_\mathcal{X})
\ar[r] \ar[d]^{Rf_*} &
D(\mathcal{O}_\mathcal{X})
\ar[d]^{Rf_*} \\
D^{+}_{\textit{Parasitic} \cap \textit{LQCoh}^{fbc}}(\mathcal{O}_\mathcal{Y})
\ar[r] &
D^{+}_{\textit{LQCoh}^{fbc}}(\mathcal{O}_\mathcal{Y}) \ar[r] &
D(\mathcal{O}_\mathcal{Y})
}
$$
et induit donc un foncteur
$$
Rf_{\QCoh, *} :
D^{+}_\QCoh(\mathcal{O}_\mathcal{X})
\longrightarrow
D^{+}_\QCoh(\mathcal{O}_\mathcal{Y})
$$
sur les catégories quotients. De plus, les foncteurs $R^if_\QCoh$ de
Cohomologie des champs, proposition
\ref{stacks-cohomology-proposition-direct-image-quasi-coherent}, sont égaux à
$H^i \circ Rf_{\QCoh, *}$, où $H^i$ est défini dans
(\ref{equation-Hi-quasi-coherent}).
\end{proposition}

\begin{proof}
Il faut montrer que $Rf_*E$ est un objet de
$D^{+}_{\textit{LQCoh}^{fbc}}(\mathcal{O}_\mathcal{Y})$ pour $E$ appartenant à
$D^{+}_{\textit{LQCoh}^{fbc}}(\mathcal{O}_\mathcal{X})$. Cela résulte de
Cohomologie des champs, proposition
\ref{stacks-cohomology-proposition-loc-qcoh-flat-base-change}, et de la suite
spectrale $R^if_*H^j(E) \Rightarrow R^{i + j}f_*E$. Le cas des modules
parasites se traite de la même manière à l'aide de Cohomologie des champs,
lemme \ref{stacks-cohomology-lemma-pushforward-parasitic}. La dernière
assertion découle immédiatement de la définition de $H^i$ dans
(\ref{equation-Hi-quasi-coherent}).
\end{proof}




\section{Image inverse dérivée des modules quasi-cohérents}
\label{section-derived-pullback}

\noindent
L'image inverse dérivée des complexes à faisceaux de cohomologie
quasi-cohérents existe en toute généralité.

\begin{proposition}
\label{proposition-derived-pullback-quasi-coherent}
Soit $f : \mathcal{X} \to \mathcal{Y}$ un morphisme de champs algébriques. Le
foncteur exact $f^*$ induit un diagramme commutatif
$$
\xymatrix{
D_{\textit{LQCoh}^{fbc}}(\mathcal{O}_\mathcal{X}) \ar[r] &
D(\mathcal{O}_\mathcal{X}) \\
D_{\textit{LQCoh}^{fbc}}(\mathcal{O}_\mathcal{Y})
\ar[r] \ar[u]^{f^*} &
D(\mathcal{O}_\mathcal{Y}) \ar[u]^{f^*}
}
$$
Le composé
$$
D_{\textit{LQCoh}^{fbc}}(\mathcal{O}_\mathcal{Y})
\xrightarrow{f^*}
D_{\textit{LQCoh}^{fbc}}(\mathcal{O}_\mathcal{X})
\xrightarrow{q_\mathcal{X}}
D_\QCoh(\mathcal{O}_\mathcal{X})
$$
est dérivable à gauche relativement à la localisation
$D_{\textit{LQCoh}^{fbc}}(\mathcal{O}_\mathcal{Y}) \to
D_\QCoh(\mathcal{O}_\mathcal{Y})$
et nous pouvons définir $Lf^*_\QCoh$ comme son foncteur dérivé gauche
$$
Lf_\QCoh^* :
D_\QCoh(\mathcal{O}_\mathcal{Y})
\longrightarrow
D_\QCoh(\mathcal{O}_\mathcal{X})
$$
(voir Catégories dérivées, définitions
\ref{derived-definition-right-derived-functor-defined} et
\ref{derived-definition-everywhere-defined}). Si $f$ est quasi-compact et
quasi-séparé, alors $Lf^*_\QCoh$ et $Rf_{\QCoh, *}$ satisfont à l'adjonction
suivante :
$$
\Hom_{D_\QCoh(\mathcal{O}_\mathcal{X})}(Lf^*_\QCoh A, B)
=
\Hom_{D_\QCoh(\mathcal{O}_\mathcal{Y})}(A, Rf_{\QCoh, *}B)
$$
pour $A \in D_\QCoh(\mathcal{O}_\mathcal{Y})$ et
$B \in D^{+}_\QCoh(\mathcal{O}_\mathcal{X})$.
\end{proposition}

\begin{proof}
Pour démontrer la première assertion, il faut montrer que $f^*E$ est un objet
de $D_{\textit{LQCoh}^{fbc}}(\mathcal{O}_\mathcal{X})$ pour $E$ appartenant à
$D_{\textit{LQCoh}^{fbc}}(\mathcal{O}_\mathcal{Y})$. Puisque
$f^* = f^{-1}$ est exact, cela résulte immédiatement du fait que $f^*$ envoie
$\textit{LQCoh}^{fbc}(\mathcal{O}_\mathcal{Y})$ dans
$\textit{LQCoh}^{fbc}(\mathcal{O}_\mathcal{X})$ d'après Cohomologie des champs,
proposition
\ref{stacks-cohomology-proposition-loc-qcoh-flat-base-change}.

\medskip\noindent
Posons $\mathcal{D} = D_{\textit{LQCoh}^{fbc}}(\mathcal{O}_\mathcal{Y})$. Soit
$S$ la collection des morphismes de $\mathcal{D}$ dont le cône est un objet de
$D_{\textit{Parasitic} \cap \textit{LQCoh}^{fbc}}(\mathcal{O}_\mathcal{Y})$.
Posons $\mathcal{D}' = D_\QCoh(\mathcal{O}_\mathcal{X})$. Posons
$F = q_\mathcal{X} \circ f^* : \mathcal{D} \to \mathcal{D}'$. Alors
$\mathcal{D}, S, \mathcal{D}', F$ sont comme dans Catégories dérivées,
situation \ref{derived-situation-derived-functor} et définition
\ref{derived-definition-right-derived-functor-defined}. Montrons que $LF(E)$
est défini pour tout objet $E$ de $\mathcal{D}$. Considérons en effet le
triangle
$$
E' \to E \to P \to E'[1]
$$
construit au lemme \ref{lemma-bousfield-colocalization}. Remarquons que
$s : E' \to E$ est un élément de $S$. Nous affirmons que $E'$ calcule $LF$.
En effet, supposons que $s' : E'' \to E$ soit un autre élément de $S$, c'est-à-dire
qu'il s'insère dans un triangle $E'' \to E \to P' \to E''[1]$ avec $P'$
appartenant à
$D_{\textit{Parasitic} \cap \textit{LQCoh}^{fbc}}(\mathcal{O}_\mathcal{Y})$.
D'après le lemme \ref{lemma-bousfield-colocalization} et sa démonstration,
$E' \to E$ se factorise par $E'' \to E$. Nous voyons ainsi que $E' \to E$ est
cofinal dans le système $S/E$. Il est donc clair que $E'$ calcule $LF$.

\medskip\noindent
Pour établir la dernière assertion, écrivons $B = q_\mathcal{X}(H)$ et
$A = q_\mathcal{Y}(E)$. Choisissons $E' \to E$ comme ci-dessus. Nous
utiliserons d'une part l'égalité
$Rf_{\QCoh, *}(B) = q_\mathcal{Y}(Rf_*H)$
et d'autre part l'égalité
$Lf^*_\QCoh(A) = q_\mathcal{X}(f^*E')$.
\begin{align*}
\Hom_{D_\QCoh(\mathcal{O}_\mathcal{X})}(Lf^*_\QCoh A, B)
& = 
\Hom_{D_\QCoh(\mathcal{O}_\mathcal{X})}(q_\mathcal{X}(f^*E'),
q_\mathcal{X}(H)) \\
& = 
\colim_{H \to H'} \Hom_{D(\mathcal{O}_\mathcal{X})}(f^*E', H') \\
& = \colim_{H \to H'} \Hom_{D(\mathcal{O}_\mathcal{Y})}(E', Rf_*H') \\
& = \Hom_{D(\mathcal{O}_\mathcal{Y})}(E', Rf_*H) \\
& =
\Hom_{D_\QCoh(\mathcal{O}_\mathcal{Y})}(A, Rf_{\QCoh, *}B)
\end{align*}
Ici, la colimite porte sur les morphismes $s : H \to H'$ de
$D^+_{\textit{LQCoh}^{fbc}}(\mathcal{O}_\mathcal{X})$
dont le cône $P(s)$ est un objet de
$D^+_{\textit{Parasitic} \cap \textit{LQCoh}^{fbc}}(\mathcal{O}_\mathcal{X})$.
Nous avons vu la première égalité ci-dessus. La deuxième résulte de la
construction du quotient de Verdier. La troisième découle de Cohomologie sur
les sites, lemme \ref{sites-cohomology-lemma-adjoint}. Comme $Rf_*P(s)$ est un
objet de
$D^+_{\textit{Parasitic} \cap \textit{LQCoh}^{fbc}}(\mathcal{O}_\mathcal{Y})$
d'après la proposition
\ref{proposition-derived-direct-image-quasi-coherent}, nous avons
$\Hom_{D(\mathcal{O}_\mathcal{Y})}(E', Rf_*P(s)) = 0$. La quatrième égalité est
donc vérifiée. La dernière résulte de la construction de $E'$.
\end{proof}








\section{Objets quasi-cohérents dans la catégorie dérivée}
\label{section-QC}

\noindent
Cette section prolonge Faisceaux sur les champs, section
\ref{stacks-sheaves-section-QC}. Soit $\mathcal{X}$ un champ algébrique. Dans
cette section, nous avons défini une catégorie triangulée
$$
\mathit{QC}(\mathcal{X}) = \mathit{QC}(\mathcal{X}_{affine}, \mathcal{O})
$$
et nous avons démontré que, si $\mathcal{X}$ est représentable par un espace
algébrique $X$, alors $\mathit{QC}(\mathcal{X})$ est équivalente à
$D_\QCoh(\mathcal{O}_X)$. Il se trouve que la théorie développée jusqu'ici
suffit exactement à démontrer le même résultat pour tout champ algébrique.

\begin{lemma}
\label{lemma-cohomology-parasitic}
Soit $\mathcal{X}$ un champ algébrique. Soit $K$ un objet de
$D(\mathcal{X}_{fppf})$ dont les faisceaux de cohomologie sont parasites. Alors
$R\Gamma(x, K) = 0$ pour tout objet $x$ de $\mathcal{X}$ situé au-dessus d'un
schéma $U$ tel que $U \to \mathcal{X}$ soit plat.
\end{lemma}

\begin{proof}
Notons $g : \Sh(\mathcal{X}_{flat, fppf}) \to \Sh(\mathcal{X}_{fppf})$ le
morphisme de topos considéré à la section \ref{section-lisse-etale}. Soit $x$
un objet de $\mathcal{X}$ situé au-dessus d'un schéma $U$ tel que
$U \to \mathcal{X}$ soit plat, c'est-à-dire un objet $x$ de
$\mathcal{X}_{flat, fppf}$. D'après la partie (2)(b) du lemme
\ref{lemma-lisse-etale-cohomology}, nous avons
$R\Gamma(x, K) = R\Gamma(\mathcal{X}_{flat, fppf}/x, g^{-1}K)$.
Or notre hypothèse signifie que les faisceaux de cohomologie de l'objet
$g^{-1}K$ de $D(\mathcal{X}_{flat, fppf})$ sont nuls ; voir Cohomologie des
champs, définition \ref{stacks-cohomology-definition-parasitic}. Ainsi,
$g^{-1}K = 0$, ce qui conclut.
\end{proof}

\begin{lemma}
\label{lemma-cohomology-parasitic-complex}
Soit $\mathcal{X}$ un champ algébrique. Soit $K$ un objet de
$D(\mathcal{X}_{fppf})$ tel que $R\Gamma(x, K) = 0$ pour tout objet $x$ de
$\mathcal{X}$ situé au-dessus d'un schéma affine $U$ tel que
$U \to \mathcal{X}$ soit plat. Alors $H^i(\mathcal{X}, K) = 0$ pour tout $i$.
\end{lemma}

\begin{proof}
Notons $g : \Sh(\mathcal{X}_{flat, fppf}) \to \Sh(\mathcal{X}_{fppf})$ le
morphisme de topos considéré à la section \ref{section-lisse-etale}. D'après la
partie (2)(b) du lemme \ref{lemma-lisse-etale-cohomology}, notre hypothèse
signifie que $g^{-1}K$ a une cohomologie nulle sur tout objet de
$\mathcal{X}_{flat, fppf}$ situé au-dessus d'un schéma affine. Comme tout objet
$x$ de $\mathcal{X}_{flat, fppf}$ possède un recouvrement par de tels objets,
nous en déduisons que les faisceaux de cohomologie de $g^{-1}K$ sont nuls,
c'est-à-dire que $g^{-1}K = 0$. Bien entendu,
$R\Gamma(\mathcal{X}_{flat, fppf}, g^{-1}K) = 0$, ce qui donne le résultat voulu
d'après la partie (2)(a) du lemme \ref{lemma-lisse-etale-cohomology}.
\end{proof}

\begin{lemma}
\label{lemma-higher-shriek-QC}
Soit $\mathcal{X}$ un champ algébrique. Soit $K$ un objet de
$D_\QCoh(\mathcal{O}_{\mathcal{X}_{flat, fppf}})$. Alors $Lg_!K$ possède la
propriété suivante : pour tout morphisme $x \to x'$ de
$\mathcal{X}_{affine}$, l'application
$$
R\Gamma(x', Lg_!K) \otimes_{\mathcal{O}(x')}^\mathbf{L} \mathcal{O}(x)
\longrightarrow
R\Gamma(x, Lg_!K)
$$
est un quasi-isomorphisme.
\end{lemma}

\begin{proof}
D'après la partie (2)(c) du lemme \ref{lemma-derived-quasi-coherent}, l'objet
$Lg_!K$ appartient à
$D_{\textit{LQCoh}^{fbc}}(\mathcal{O}_\mathcal{X})$. Il en résulte aussitôt que
l'application affichée dans le lemme est un isomorphisme lorsque
$\mathcal{O}(x') \to \mathcal{O}(x)$ est un homomorphisme d'anneaux plat ; nous
omettons les détails.

\medskip\noindent
Montrons dans ce paragraphe que la question est locale pour la topologie
étale. Soit $x \to x'$ un morphisme arbitraire de $\mathcal{X}_{affine}$. Soit
$\{x'_i \to x'\}$ un recouvrement dans
$\mathcal{X}_{affine, \etale}$. Posons $x_i = x \times_{x'} x'_i$, de sorte que
$\{x_i \to x\}$ soit aussi un recouvrement de
$\mathcal{X}_{affine, \etale}$. Alors
$\mathcal{O}(x') \to \prod \mathcal{O}(x'_i)$ est un homomorphisme d'anneaux
étale fidèlement plat et
$$
\prod \mathcal{O}(x_i) =
\mathcal{O}(x) \otimes_{\mathcal{O}(x')} \left(\prod \mathcal{O}(x'_i)\right)
$$
Un argument élémentaire d'algèbre, que nous omettons, montre donc qu'il suffit
d'établir le résultat de l'énoncé pour chacun des morphismes
$x_i \to x'_i$ de $\mathcal{X}_{affine}$. Autrement dit, le problème est local
pour la topologie étale.

\medskip\noindent
Choisissons un schéma $X$ et un morphisme lisse surjectif
$f : X \to \mathcal{X}$. Nous pouvons considérer $f$ comme un objet de
$\mathcal{X}$, conformément à notre abus de notation ; alors
$(\Sch/X)_{fppf} = \mathcal{X}/f$ ; voir Faisceaux sur les champs, section
\ref{stacks-sheaves-section-restriction}. D'après Faisceaux sur les champs,
lemme
\ref{stacks-sheaves-lemma-surjective-flat-locally-finite-presentation}
par exemple, il existe un recouvrement étale $\{x'_i \to x'\}$ tel que
$x'_i : U'_i = p(x'_i) \to \mathcal{X}$ se factorise par $f$. D'après le
résultat du paragraphe précédent, nous pouvons supposer que $x \to x'$ est
l'image d'un morphisme $U \to U'$ de $(\textit{Aff}/X)_{fppf}$ par le foncteur
$(\Sch/X)_{fppf} \to \mathcal{X}$. À ce stade, nous utilisons le fait que la
restriction à $(\Sch/X)_{fppf}$ de $Lg_!K$ est égale à
$f^*Lg_!K = L(g')_!(f')^*K$ d'après le lemme
\ref{lemma-lisse-etale-functorial-derived}. Nous sommes ainsi ramenés au cas
étudié au paragraphe suivant.

\medskip\noindent
Supposons $\mathcal{X} = (\Sch/X)_{fppf}$ et que $x \to x'$ corresponde au
morphisme de schémas affines $U \to U'$. Nous pouvons encore travailler
localement pour la topologie étale, ou de Zariski, sur $U'$ ; nous pouvons donc
supposer que $U' \to X$ se factorise par un ouvert affine de $X$. Nous sommes
ainsi ramenés au cas du paragraphe suivant.

\medskip\noindent
Supposons $\mathcal{X} = (\Sch/X)_{fppf}$, où $X = \Spec(R)$ est un schéma
affine, et que $x \to x'$ corresponde au morphisme de schémas affines
$U \to U'$. Soit $M^\bullet$ un complexe de $R$-modules représentant
$R\Gamma(X, K)$. La construction de Compléments d'algèbre, lemme
\ref{more-algebra-lemma-K-flat-resolution}, permet de supposer que
$M^\bullet = \colim P_n^\bullet$, où chaque $P_n^\bullet$ est un complexe borné
supérieurement de $R$-modules libres. Nous omettons les détails ; voir aussi
Compléments d'algèbre, remarque \ref{more-algebra-remark-P-resolution}.
Considérons le complexe de modules $M^\bullet_{flat, fppf}$ sur
$X_{flat, fppf} = (\Sch/X)_{flat, fppf}$ défini par la règle
$$
U \longmapsto \Gamma(U, M^\bullet \otimes_R \mathcal{O}_U)
$$
C'est un complexe de faisceaux d'après la discussion de Descente, section
\ref{descent-section-quasi-coherent-sheaves}. Il existe une application
canonique $M^\bullet_{flat, fppf} \to K$ qui, d'après les remarques faites au
début de la démonstration, induit un isomorphisme sur les sections au-dessus des
objets affines de $X_{flat, fppf}$. Comme tout objet de $X_{flat, fppf}$ possède
un recouvrement par des objets affines, nous voyons que
$M^\bullet_{flat, fppf}$ coïncide avec $K$.

\medskip\noindent
Soit $M^\bullet_{fppf}$ le complexe de modules sur $X_{fppf}$ donné par la
même formule que ci-dessus. Rappelons que
$Lg_!\mathcal{O} = g_!\mathcal{O} = \mathcal{O}$. Puisque $Lg_!$ est le
foncteur dérivé gauche de $g_!$, nous en déduisons que
$Lg_!P_{n, flat, fppf}^\bullet = P_{n, fppf}^\bullet$. Comme le foncteur $Lg_!$
commute aux colimites homotopiques, ou directement par sa construction dans
Cohomologie sur les sites, lemme
\ref{sites-cohomology-lemma-existence-derived-lower-shriek}, et comme
$M^\bullet = \colim P_n^\bullet$, nous concluons que
$Lg_!M^\bullet_{flat, fppf} = M^\bullet_{fppf}$. Écrivons
$U = \Spec(A)$, $U' = \Spec(A')$ et supposons que $U \to U'$ corresponde à
l'homomorphisme d'anneaux $A' \to A$. D'après ce qui précède,
$$
R\Gamma(U, Lg_!K) = M^\bullet \otimes_R A
\quad\text{et}\quad
R\Gamma(U', Lg_!K) = M^\bullet \otimes_R A'
$$
Puisque $M^\bullet$ est un complexe K-plat de $R$-modules, la transitivité du
produit tensoriel montre que
$$
R\Gamma(U', Lg_!K) \otimes_{A'}^\mathbf{L} A
\longrightarrow
R\Gamma(U, Lg_!K)
$$
est le quasi-isomorphisme recherché.
\end{proof}

\begin{proposition}
\label{proposition-QC-compare}
Soit $\mathcal{X}$ un champ algébrique. Alors $\mathit{QC}(\mathcal{X})$ est
canoniquement équivalente à $D_\QCoh(\mathcal{O}_\mathcal{X})$.
\end{proposition}

\begin{proof}
D'après Faisceaux sur les champs, lemme
\ref{stacks-sheaves-lemma-QC-compare-fppf}, l'image inverse par le morphisme de
comparaison
$\epsilon : \mathcal{X}_{affine, fppf} \to \mathcal{X}_{affine}$
identifie $\mathit{QC}(\mathcal{X})$ à une sous-catégorie pleine
$Q_\mathcal{X} \subset D(\mathcal{X}_{affine, fppf}, \mathcal{O})$. À l'aide
de l'équivalence de topos annelés de Faisceaux sur les champs, équation
(\ref{stacks-sheaves-equation-alternative-ringed}), nous pouvons et allons
considérer $Q_\mathcal{X}$ comme une sous-catégorie pleine de
$D(\mathcal{X}_{fppf}, \mathcal{O})$.

\medskip\noindent
De même, le lemme \ref{lemma-bousfield-colocalization} et la remarque
\ref{remark-QCoh-admissible} montrent que
$D_\QCoh(\mathcal{O}_\mathcal{X})$ peut être considérée comme l'orthogonal
gauche $\mathcal{A}$ de la sous-catégorie admissible à gauche
$D_{\textit{Parasitic} \cap \textit{LQCoh}^{fbc}}(\mathcal{O}_\mathcal{X})$
de $D_{\textit{LQCoh}^{fbc}}(\mathcal{O}_\mathcal{X})$.

\medskip\noindent
Pour conclure, nous montrerons que $Q_\mathcal{X}$ est égale à $\mathcal{A}$
comme sous-catégorie de $D(\mathcal{X}_{fppf}, \mathcal{O})$.

\medskip\noindent
Étape 1 : $Q_\mathcal{X}$ est contenue dans
$D_{\textit{LQCoh}^{fbc}}(\mathcal{O}_\mathcal{X})$.
Un objet $K$ de $Q_\mathcal{X}$ est caractérisé par la propriété suivante :
$K$, considéré comme un objet de
$D(\mathcal{X}_{affine, fppf}, \mathcal{O})$, est tel que $R\epsilon_*K$ soit un
objet de $\mathit{QC}(\mathcal{X}_{affine}, \mathcal{O})$. Cela signifie
exactement que, pour tout morphisme $x \to x'$ de
$\mathcal{X}_{affine}$, l'application
$$
R\Gamma(x', K) \otimes_{\mathcal{O}(x')}^\mathbf{L} \mathcal{O}(x)
\longrightarrow
R\Gamma(x, K)
$$
est un isomorphisme ; voir la note de bas de page dans l'énoncé de Cohomologie
sur les sites, lemme
\ref{sites-cohomology-lemma-cartesion-plus-topology}. Maintenant, si
$x' \to x$ est situé au-dessus d'un morphisme plat de schémas affines, cela
signifie que
$$
H^i(x', K) \otimes_{\mathcal{O}(x')} \mathcal{O}(x)
\cong
H^i(x, K)
$$
Cela signifie clairement que $H^i(K)$ est un faisceau pour la topologie étale
(Faisceaux sur les champs, lemme
\ref{stacks-sheaves-lemma-quasi-coherent-alternative}) et qu'il possède la
propriété de changement de base plat ; nous omettons un petit détail.

\medskip\noindent
Étape 2 : $Q_\mathcal{X}$ est contenue dans $\mathcal{A}$. Il suffit pour cela
de montrer que, pour $K$ dans $Q_\mathcal{X}$, nous avons
$\Hom(K, P) = 0$ pour tout $P$ dans
$D_{\textit{Parasitic} \cap \textit{LQCoh}^{fbc}}(\mathcal{O}_\mathcal{X})$.
Considérons l'objet
$$
H = R\SheafHom_{\mathcal{O}_\mathcal{X}}(K, P)
$$
Soit $x$ un objet de $\mathcal{X}$ situé au-dessus d'un schéma affine
$U = p(x)$. D'après Cohomologie sur les sites, lemme
\ref{sites-cohomology-lemma-section-RHom-over-U}, nous avons la première
égalité de
$$
R\Gamma(x, H) =
R\Hom_{\mathcal{O}_\mathcal{X}}(K|_{\mathcal{X}/x}, P|_{\mathcal{X}/x}) =
R\Hom_{\mathcal{O}}(K|_{\mathcal{X}_{affine}/x}, P|_{\mathcal{X}_{affine}/x})
$$
La seconde égalité vient de ce que le topos du site $\mathcal{X}/x$ est
équivalent au topos du site $\mathcal{X}_{affine}/x$ ; voir Faisceaux sur les
champs, équation (\ref{stacks-sheaves-equation-alternative-ringed}). Nous
pouvons écrire $K = \epsilon^*N$ pour un certain $N$ de
$\mathit{QC}(\mathcal{O})$. Alors Cohomologie sur les sites, lemme
\ref{sites-cohomology-lemma-QC-hom-out-of-plus-topology}, donne
$$
R\Gamma(x, H) =
R\Hom_{D(\mathcal{O}(x))}(R\Gamma(x, N), R\Gamma(x, P))
$$
Le lemme \ref{lemma-cohomology-parasitic} montre que
$R\Gamma(x, P) = 0$ si $U \to \mathcal{X}$ est plat, et donc que
$R\Gamma(x, H) = 0$ sous la même hypothèse. D'après le lemme
\ref{lemma-cohomology-parasitic-complex}, nous concluons que
$R\Gamma(\mathcal{X}, H) = 0$, et par conséquent que $\Hom(K, P) = 0$.

\medskip\noindent
Étape 3 : $\mathcal{A}$ est contenue dans $Q_\mathcal{X}$. Soit $K$ un objet de
$\mathcal{A}$ et soit $x \to x'$ un morphisme de
$\mathcal{X}_{affine}$. Il faut montrer que
$$
R\Gamma(x', K) \otimes_{\mathcal{O}(x')}^\mathbf{L} \mathcal{O}(x)
\longrightarrow
R\Gamma(x, K)
$$
est un quasi-isomorphisme ; voir la note de bas de page dans l'énoncé de
Cohomologie sur les sites, lemme
\ref{sites-cohomology-lemma-cartesion-plus-topology}. La démonstration du lemme
\ref{lemma-bousfield-colocalization} et la discussion de la remarque
\ref{remark-QCoh-admissible} montrent que $\mathcal{A}$ est l'image de la
restriction de $Lg_!$ à
$D_\QCoh(\mathcal{O}_{\mathcal{X}_{flat, fppf}})$. Nous pouvons donc supposer
$K = Lg_!M$ pour un certain $M$ de
$D_\QCoh(\mathcal{O}_{\mathcal{X}_{flat, fppf}})$. L'égalité recherchée résulte
alors du lemme \ref{lemma-higher-shriek-QC}.
\end{proof}


















\begin{multicols}{2}[\section{Autres chapitres}]
\noindent
Préliminaires
\begin{enumerate}
\item \hyperref[introduction-section-phantom]{Introduction}
\item \hyperref[conventions-section-phantom]{Conventions}
\item \hyperref[sets-section-phantom]{Théorie des ensembles}
\item \hyperref[categories-section-phantom]{Catégories}
\item \hyperref[topology-section-phantom]{Topologie}
\item \hyperref[sheaves-section-phantom]{Faisceaux sur les espaces}
\item \hyperref[sites-section-phantom]{Sites et faisceaux}
\item \hyperref[stacks-section-phantom]{Champs}
\item \hyperref[fields-section-phantom]{Corps}
\item \hyperref[algebra-section-phantom]{Algèbre commutative}
\item \hyperref[brauer-section-phantom]{Groupes de Brauer}
\item \hyperref[homology-section-phantom]{Algèbre homologique}
\item \hyperref[derived-section-phantom]{Catégories dérivées}
\item \hyperref[simplicial-section-phantom]{Méthodes simpliciales}
\item \hyperref[more-algebra-section-phantom]{Compléments d'algèbre}
\item \hyperref[smoothing-section-phantom]{Lissage des morphismes d'anneaux}
\item \hyperref[modules-section-phantom]{Faisceaux de modules}
\item \hyperref[sites-modules-section-phantom]{Modules sur les sites}
\item \hyperref[injectives-section-phantom]{Objets injectifs}
\item \hyperref[cohomology-section-phantom]{Cohomologie des faisceaux}
\item \hyperref[sites-cohomology-section-phantom]{Cohomologie sur les sites}
\item \hyperref[dga-section-phantom]{Algèbre différentielle graduée}
\item \hyperref[dpa-section-phantom]{Algèbre à puissances divisées}
\item \hyperref[sdga-section-phantom]{Faisceaux différentiels gradués}
\item \hyperref[hypercovering-section-phantom]{Hyperrecouvrements}
\end{enumerate}
Schémas
\begin{enumerate}
\setcounter{enumi}{25}
\item \hyperref[schemes-section-phantom]{Schémas}
\item \hyperref[constructions-section-phantom]{Constructions de schémas}
\item \hyperref[properties-section-phantom]{Propriétés des schémas}
\item \hyperref[morphisms-section-phantom]{Morphismes de schémas}
\item \hyperref[coherent-section-phantom]{Cohomologie des schémas}
\item \hyperref[divisors-section-phantom]{Diviseurs}
\item \hyperref[limits-section-phantom]{Limites de schémas}
\item \hyperref[varieties-section-phantom]{Variétés}
\item \hyperref[topologies-section-phantom]{Topologies sur les schémas}
\item \hyperref[descent-section-phantom]{Descente}
\item \hyperref[perfect-section-phantom]{Catégories dérivées de schémas}
\item \hyperref[more-morphisms-section-phantom]{Compléments sur les morphismes}
\item \hyperref[flat-section-phantom]{Compléments sur la platitude}
\item \hyperref[groupoids-section-phantom]{Schémas en groupoïdes}
\item \hyperref[more-groupoids-section-phantom]{Compléments sur les schémas en groupoïdes}
\item \hyperref[etale-section-phantom]{Morphismes étales de schémas}
\end{enumerate}
Thèmes de théorie des schémas
\begin{enumerate}
\setcounter{enumi}{41}
\item \hyperref[chow-section-phantom]{Homologie de Chow}
\item \hyperref[intersection-section-phantom]{Théorie de l'intersection}
\item \hyperref[pic-section-phantom]{Schémas de Picard des courbes}
\item \hyperref[weil-section-phantom]{Théories cohomologiques de Weil}
\item \hyperref[adequate-section-phantom]{Modules adéquats}
\item \hyperref[dualizing-section-phantom]{Complexes dualisants}
\item \hyperref[duality-section-phantom]{Dualité pour les schémas}
\item \hyperref[discriminant-section-phantom]{Discriminants et différentes}
\item \hyperref[derham-section-phantom]{Cohomologie de de Rham}
\item \hyperref[local-cohomology-section-phantom]{Cohomologie locale}
\item \hyperref[algebraization-section-phantom]{Géométrie algébrique et formelle}
\item \hyperref[curves-section-phantom]{Courbes algébriques}
\item \hyperref[resolve-section-phantom]{Résolution des surfaces}
\item \hyperref[models-section-phantom]{Réduction semi-stable}
\item \hyperref[functors-section-phantom]{Foncteurs et morphismes}
\item \hyperref[equiv-section-phantom]{Catégories dérivées de variétés}
\item \hyperref[pione-section-phantom]{Groupes fondamentaux de schémas}
\item \hyperref[etale-cohomology-section-phantom]{Cohomologie étale}
\item \hyperref[crystalline-section-phantom]{Cohomologie cristalline}
\item \hyperref[proetale-section-phantom]{Cohomologie pro-étale}
\item \hyperref[relative-cycles-section-phantom]{Cycles relatifs}
\item \hyperref[more-etale-section-phantom]{Compléments de cohomologie étale}
\item \hyperref[trace-section-phantom]{La formule des traces}
\end{enumerate}
Espaces algébriques
\begin{enumerate}
\setcounter{enumi}{64}
\item \hyperref[spaces-section-phantom]{Espaces algébriques}
\item \hyperref[spaces-properties-section-phantom]{Propriétés des espaces algébriques}
\item \hyperref[spaces-morphisms-section-phantom]{Morphismes d'espaces algébriques}
\item \hyperref[decent-spaces-section-phantom]{Espaces algébriques décents}
\item \hyperref[spaces-cohomology-section-phantom]{Cohomologie des espaces algébriques}
\item \hyperref[spaces-limits-section-phantom]{Limites d'espaces algébriques}
\item \hyperref[spaces-divisors-section-phantom]{Diviseurs sur les espaces algébriques}
\item \hyperref[spaces-over-fields-section-phantom]{Espaces algébriques sur des corps}
\item \hyperref[spaces-topologies-section-phantom]{Topologies sur les espaces algébriques}
\item \hyperref[spaces-descent-section-phantom]{Descente et espaces algébriques}
\item \hyperref[spaces-perfect-section-phantom]{Catégories dérivées d'espaces}
\item \hyperref[spaces-more-morphisms-section-phantom]{Compléments sur les morphismes d'espaces}
\item \hyperref[spaces-flat-section-phantom]{Platitude sur les espaces algébriques}
\item \hyperref[spaces-groupoids-section-phantom]{Groupoïdes dans les espaces algébriques}
\item \hyperref[spaces-more-groupoids-section-phantom]{Compléments sur les groupoïdes dans les espaces}
\item \hyperref[bootstrap-section-phantom]{Amorçage}
\item \hyperref[spaces-pushouts-section-phantom]{Sommes amalgamées d'espaces algébriques}
\end{enumerate}
Thèmes de géométrie
\begin{enumerate}
\setcounter{enumi}{81}
\item \hyperref[spaces-chow-section-phantom]{Groupes de Chow des espaces}
\item \hyperref[groupoids-quotients-section-phantom]{Quotients de groupoïdes}
\item \hyperref[spaces-more-cohomology-section-phantom]{Compléments sur la cohomologie des espaces}
\item \hyperref[spaces-simplicial-section-phantom]{Espaces simpliciaux}
\item \hyperref[spaces-duality-section-phantom]{Dualité pour les espaces}
\item \hyperref[formal-spaces-section-phantom]{Espaces algébriques formels}
\item \hyperref[restricted-section-phantom]{Algébrisation des espaces formels}
\item \hyperref[spaces-resolve-section-phantom]{Résolution des surfaces revisitée}
\end{enumerate}
Théorie des déformations
\begin{enumerate}
\setcounter{enumi}{89}
\item \hyperref[formal-defos-section-phantom]{Théorie formelle des déformations}
\item \hyperref[defos-section-phantom]{Théorie des déformations}
\item \hyperref[cotangent-section-phantom]{Le complexe cotangent}
\item \hyperref[examples-defos-section-phantom]{Problèmes de déformation}
\end{enumerate}
Champs algébriques
\begin{enumerate}
\setcounter{enumi}{93}
\item \hyperref[algebraic-section-phantom]{Champs algébriques}
\item \hyperref[examples-stacks-section-phantom]{Exemples de champs}
\item \hyperref[stacks-sheaves-section-phantom]{Faisceaux sur les champs algébriques}
\item \hyperref[criteria-section-phantom]{Critères de représentabilité}
\item \hyperref[artin-section-phantom]{Axiomes d'Artin}
\item \hyperref[quot-section-phantom]{Espaces de Quot et de Hilbert}
\item \hyperref[stacks-properties-section-phantom]{Propriétés des champs algébriques}
\item \hyperref[stacks-morphisms-section-phantom]{Morphismes de champs algébriques}
\item \hyperref[stacks-limits-section-phantom]{Limites de champs algébriques}
\item \hyperref[stacks-cohomology-section-phantom]{Cohomologie des champs algébriques}
\item \hyperref[stacks-perfect-section-phantom]{Catégories dérivées de champs}
\item \hyperref[stacks-introduction-section-phantom]{Introduction aux champs algébriques}
\item \hyperref[stacks-more-morphisms-section-phantom]{Compléments sur les morphismes de champs}
\item \hyperref[stacks-geometry-section-phantom]{La géométrie des champs}
\end{enumerate}
Thèmes de théorie des espaces de modules
\begin{enumerate}
\setcounter{enumi}{107}
\item \hyperref[moduli-section-phantom]{Champs de modules}
\item \hyperref[moduli-curves-section-phantom]{Espaces de modules de courbes}
\end{enumerate}
Divers
\begin{enumerate}
\setcounter{enumi}{109}
\item \hyperref[examples-section-phantom]{Exemples}
\item \hyperref[exercises-section-phantom]{Exercices}
\item \hyperref[guide-section-phantom]{Guide bibliographique}
\item \hyperref[desirables-section-phantom]{Desiderata}
\item \hyperref[coding-section-phantom]{Style de programmation}
\item \hyperref[obsolete-section-phantom]{Obsolète}
\item \hyperref[fdl-section-phantom]{Licence de documentation libre GNU}
\item \hyperref[index-section-phantom]{Index généré automatiquement}
\end{enumerate}
\end{multicols}


\bibliography{my}
\bibliographystyle{amsalpha}

\end{document}
